\section{Setup}
\label{sec:setup}

An episode has three stages (Figure~\ref{fig:setup}a). Each of two copies
of a model receives a private card with its assignment and writes a short
note on its plan. Both then write a 48-token reflection, during which
A's attention also reads B's memory through the memory link. Finally,
from the state saved after the reflection, A answers questions about
its own assignment, B's, and Robin's, either with the link kept or after
it has been cut.

\paragraph{Private cards.}
Testing attribution needs content whose owner we know, so we write that
content ourselves. The two copies, the receiver A and the partner B, are
instances of one frozen model, Qwen3-4B-Instruct-2507
\citep{qwen2025qwen3_4b_instruct_2507,yang2025qwen3}. Both cards pose
the same choice among four delivery services, and each gives its holder
an assignment: a decision rule, stated as a priority (lowest cost, fastest
delivery, highest reliability, or lowest emissions), and a code word
such as \emph{queen}. Both cards also describe, in identical words, the
assignment of Robin, a colleague who takes no part. A, B, and Robin hold
different rules and different code words. Each card names its holder
and its partner, so the memory that A reads from B says in words whose
it is: B's card reads \qq{You are participant \{B\}. Another
participant, \{A\}, is working on this same decision in parallel\ldots},
with the participants' names in place of \{B\} and \{A\}
(Appendix~\ref{app:methods} gives the full cards and the
counterbalancing).

\paragraph{Questions.}
We then ask A which rule it was assigned, which rule it is using
now, and which code word it was assigned (the questions, like the cards,
say \qq{priority}; Appendix Table~\ref{tab:stems} gives the wording).
The assigned rule and the code word have one right answer however A's plans change, so naming B's rule as the one A was
assigned, or B's code word as A's own, is a misattribution, which we
call \emph{claiming}. Reading a partner may legitimately change the rule
A is using now, so naming B's rule there counts as \emph{adoption}. Two
further questions ask what A thinks B was assigned, and the share of
these answers that name B's item measures \emph{access}
(Appendix~\ref{app:stats} gives the corresponding logit score). The same
questions about Robin, whose entry is identical on both cards, serve as
a control: if reading B simply made B's rule the model's favored answer
to any such question, through priming or salience, the answers about
Robin would move toward it too. Because the two cards agree about Robin,
this control does not test how A treats a third party whose entries
conflict.

\paragraph{The memory link.}
Through the memory link, A's attention at every layer reads B's memory,
the fixed key--value cache of B's card, note, and reflection prompt,
alongside A's own cache. B's keys keep the rotary positions they had in
B (the position code that the model builds into every key), so that, seen from A's queries, B's card sits at nearly the same
positions as A's own card, which follows the same template. We add
$\log w$ to A's attention logits on B's keys, so that each head's output
is an exact mixture,
\begin{equation}
o=o_{\mathrm{own}}+\beta(o_B-o_{\mathrm{own}}),\qquad
\beta=\frac{w Z_B}{Z_{\mathrm{own}}+w Z_B},
\label{eq:bridge}
\end{equation}
where $o_{\mathrm{own}}$ and $o_B$ are the head's outputs over A's own
keys alone and over B's keys alone, and $Z$ sums the exponentiated,
unweighted attention scores over the indicated keys. By choosing the
link weight $w$, we set how much of each head's output comes from B, and
nothing in the mixed output marks which part came from whom; the
account of Section~\ref{sec:account} builds on this. At $w=0$ there is
no link. At
$w=1$, B's keys enter attention on the same footing as A's own, but
equal footing is not an equal split: B's mean share of A's attention
during the questions was 0.360 (Appendix~\ref{app:evidence}). Our main
setting, $w=2$, doubles the unnormalized weight of B's keys; their share
was then 0.500.

We laid B's memory over A's own positions on purpose: it is the closest
analogue to the opening image and the hardest case for keeping sources
apart, because position then gives no cue to which memory is which. Systems such as
LatentMAS instead prepend the partner's cache \citep{zou2026latentmas},
an arrangement we did not test (Section~\ref{sec:account}).

\paperfigure{fig2_setup}{2.70in}{
\textbf{A reads B's memory while both copies reflect, then answers
questions from one saved state.}
(a) While both copies reflect, A's attention also reads B's memory;
asked which rule it was assigned, A then names B's.
(b) Each question branches from the state saved after the reflection,
with the link kept or cut.
}{fig:setup}

\paragraph{Other routes.}
If claiming depended only on how much of B's content reaches A, any
route that gave A the same access would produce the same claiming. We
test this with two other routes; a third, which adds B's hidden states
to A's, was tested in a separate study (Appendix~\ref{app:v1}). The \emph{steering vector}, a
mean-difference direction for B's rule
\citep{turner2023activation,zou2023representation}, is added to A's
residual stream at one middle layer while A reflects and answers. It
pushes A's internal state toward B's rule but gives A's
attention no memory of B to read. We set its gain to match the memory
link at $w=1$ in mean access to B's rule, because no gain in our grid
matched the stronger link at $w=2$ (Appendix~\ref{app:stats}). Matching
access holds one measure of arrival fixed, how often A names B's rule
when asked about B, while the route changes.

With the \emph{text message} we ask whether A, reading B's content in
full, keeps it apart from its own when a header names the source. We
place B's reflection from an unconnected run in A's context under one of
four headers: \qq{B shared these thoughts} (tagged), \qq{Here are some
thoughts} (untagged), \qq{Here are your own earlier thoughts}, or \qq{A
stranger shared these thoughts}. Text is not matched to the memory link
in access, and because a message may legitimately change what A plans to
do, we tested the headers on adoption.

\paragraph{Link kept or cut.}
Claiming could arise at two moments: while A answers, from reading B,
or earlier, from what A wrote while connected. What A writes while
connected is its reflection, which answers the request to \qq{reflect on
your decision and on what matters to you, and mention your code word
once}, so it states, in A's own words, what A takes its assignment to
be. To separate the two moments, A answers each question on two branches
from the state saved after the reflection. With the \emph{link kept}, A
answers while still reading B. With the \emph{link cut}, reading stops
before the question, and A answers from its own cache, which still holds
the reflection; A's keys and values for it were computed while A
attended to B. Because both branches start from the same state, their
difference isolates the effect of reading at answer time, and whatever
survives the cut is what A carried into its own memory while connected.

\paragraph{Measurement.}
Near 0 and 100\%, choice rates say little about how strongly the model
prefers one answer to another, so our tests use a source score built
from answer logits,
\begin{equation}
M=\underbrace{\bigl(z_B^{\mathrm{self}}-z_A^{\mathrm{self}}\bigr)}_{\text{A asked for its assigned rule}}
-\underbrace{\bigl(z_B^{\mathrm{Robin}}-z_R^{\mathrm{Robin}}\bigr)}_{\text{same question about Robin}},
\label{eq:metric}
\end{equation}
where $z_A$, $z_B$, and $z_R$ are the logits of the answers naming A's,
B's, and Robin's rules, so differences are log probability ratios in
nats. The first term measures how far A prefers B's rule to its own
when asked about itself, and the Robin term removes any pull toward B's
rule that is not about the self. $M_{\mathrm{NOW}}$ is the same
score on the question about A's current rule. We write $\Delta M$ for
the difference in $M$ between two conditions, the link and no link
unless we name another contrast.

\paragraph{Preregistration and samples.}
We chose every setting in pilots on separate episodes without looking
at claiming: $w=2$ gave the largest increase in access among the link
weights that kept accuracy on common-knowledge questions within five
points of no link and left the answer format intact
(Appendix~\ref{app:stats}). We then fixed each confirmatory test and
its analysis, with a predicted direction where the protocol stated one,
in a frozen protocol before drawing the fresh episodes it ran on: 600
per sample in the memory-link studies and 300 in the separate
hidden-state study. Results reported without
qualification are prespecified in this sense; we say in words when a
result is post hoc or comes from a pilot. The confirmatory studies ran in
sequence, each on fresh episodes. The \emph{first confirmatory sample}
tested claiming and the text messages. The \emph{main sample}, drawn
afterwards, tested the questions those results raised, with the matched
steering vector, the cut link, and a third-person record of B's memory
(Section~\ref{sec:third}); because it
reran the memory link at $w=2$ as their reference, most rates reported
here come from it (Appendix Table~\ref{tab:study-map}). The episode
is the statistical unit: tests use Holm correction within families and
95\% episode-bootstrap intervals, and the appendix reports every
prespecified test, including those that failed or reversed
(Appendix Tables~\ref{tab:contrasts}, \ref{tab:pair-gate},
and~\ref{tab:v1contrasts}). The code, the three protocols with their
freeze manifests, and a table that maps every figure and table to the
runs it was computed from are available at
\url{https://github.com/BeibaiDraco/mind-connecting} (Appendix~\ref{app:repro}).
